\documentclass[11pt]{article}
\usepackage[T1]{fontenc}
\usepackage[utf8]{inputenc}
\usepackage{lmodern}
\usepackage{amsmath,amssymb,amsthm,mathtools}
\usepackage[margin=1.08in]{geometry}
\usepackage{microtype}
\usepackage{xurl}
\usepackage{xcolor}
\usepackage[colorlinks=true,linkcolor=blue!45!black,citecolor=blue!45!black,urlcolor=blue!45!black]{hyperref}
\hypersetup{pdftitle={An integral degree-three Gersten counterexample},pdfauthor={OpenAI}}
\newtheorem{theorem}{Theorem}[section]
\newtheorem{lemma}[theorem]{Lemma}
\newtheorem{proposition}[theorem]{Proposition}
\theoremstyle{remark}
\newtheorem{remark}[theorem]{Remark}
\DeclareMathOperator{\Spec}{Spec}
\DeclareMathOperator{\Frac}{Frac}
\DeclareMathOperator{\Tr}{Tr}
\newcommand{\Z}{\mathbb Z}
\newcommand{\Q}{\mathbb Q}
\newcommand{\F}{\mathbb F}
\newcommand{\OO}{\mathcal O}
\newcommand{\X}{\mathcal X}
\title{An integral degree-three Gersten counterexample}
\author{OpenAI}
\date{September 26, 2026}
\begin{document}
\maketitle
\begin{abstract}
We construct a two-dimensional ramified regular local ring $A$ in mixed
characteristic $(0,5)$ for which the integral map
$K_3(A)\to K_3(\Frac A)$ has nonzero kernel. This gives a negative
answer to the unrestricted integral Gersten conjecture.
\end{abstract}
\section{Introduction}\label{sec:introduction}

For a regular local domain $B$ with fraction field $F_B$, the integral
Gersten conjecture asks whether, for each $n\geq0$, the augmented
complex of Quillen groups
\[
 0\longrightarrow K_n(B)\longrightarrow K_n(F_B)
 \longrightarrow
 \bigoplus_{\operatorname{ht}\mathfrak p=1}
 K_{n-1}(\kappa(\mathfrak p))\longrightarrow\cdots
\]
is exact. Here $\kappa(\mathfrak p)$ is the residue field at a prime,
and the later maps are the codimension-residue maps. In particular,
exactness requires the first map $K_n(B)\to K_n(F_B)$ to be injective.
We construct a ring for which this requirement fails in degree three.

Fix a primitive fifth root of unity $\zeta$, and define
\[
 L=\Q_5(\zeta),\qquad V=\OO_L,\qquad \pi=1-\zeta.
\]
Let $\X\subset\mathbb P^2_V$ be the projective hypersurface
\[
 (X-Z)^5-XY^4+Y^5+\zeta Z^5=0.
\]
On its special fiber take $m=[0:0:1]$, and set
$A=\OO_{\X,m}$. We use $G_i$ for coherent $K$-theory and
$G_i(-;\Z/5)$ for its mod-five coefficient groups.

\begin{theorem}\label{thm:main}
The ring $A$ is a two-dimensional Noetherian regular local domain
of mixed characteristic $(0,5)$, with $\pi\in\mathfrak m_A^2$.
There is an integral class $c\in K_3(A)$ such that
\[
 c|_{\Frac A}=0,
 \qquad
 c|_{A[1/\pi]}\bmod5\ne0
 \quad\text{in }G_3(A[1/\pi];\Z/5).
\]
In particular, $K_3(A)\to K_3(\Frac A)$ is not injective.
\end{theorem}

Thus the unrestricted integral Gersten complex can fail already at
its initial injection. The class $c$ is not in $5K_3(A)$: any such
multiple would have zero coefficient reduction, and hence zero
reduction after restriction to $A[1/\pi]$. This conclusion does not
determine the order of $c$ or its image in $K_3(A)\otimes\Q$.

\subsection*{Context and earlier results}

Gersten asked, among other questions, whether the transfer from the
residue field to a discrete valuation ring vanishes
\cite[Problem 9]{GerstenProblems}. Localization connects this question
to injectivity of restriction to the fraction field: the image of
the residue-field transfer is its kernel. Quillen formulated the
general regular-local conjecture and proved it for regular
semilocalizations of finite-type algebras over a field
\cite[\S7, Conjecture 5.10 and Theorem 5.11]{QuillenHigher}.
Sherman established the equicharacteristic DVR case
\cite[p.~498]{Sherman}, and Panin proved the full equicharacteristic
case \cite[Theorem A]{Panin}. These results place the remaining
integral question in mixed characteristic.

In mixed characteristic, smoothness over a DVR and the choice of
coefficients lead to different positive results. Gillet and Levine's
relative theory reduces the smooth-over-DVR setting to the DVR
case \cite{GilletLevine}. Geisser and Levine proved finite-coefficient
Gersten exactness for semilocal rings essentially smooth over a DVR,
including coefficients of residue-characteristic order
\cite[Theorem 8.2]{GeisserLevine}. Druzhinin's preprint states exactness
for essentially smooth local Henselian schemes over a DVR
\cite[Theorem 4.5]{Druzhinin}. A separate result of Skalit gives
finite-coefficient exactness for every unramified regular local
ring \cite[Corollary 5.7]{Skalit}; in mixed characteristic $(0,p)$,
unramified here means $p\in\mathfrak m\setminus\mathfrak m^2$.
His integral result is a reduction to DVRs essentially smooth over
$\Z$ \cite[Theorem 5.4]{Skalit}. Our ring satisfies
$5\in\mathfrak m_A^2$, and its special fiber over the chosen $V$
is singular at $m$; it lies outside the unramified and
smooth-over-$V$ settings just described.

The distinction between integral and coefficient classes governs the
proof. A finite-coefficient kernel class may have a nonzero coefficient
boundary and hence fail to lift integrally. Even when it lifts, its
integral generic restriction may be nonzero but divisible by the
coefficient prime. Feld's mod-three degree-two example and its integral
divisibility corollary illustrate this second issue
\cite[Theorem 1.1 and Corollary 5.1]{Feld}.

Feld's construction also provides a methodological antecedent:
push a Bott-type class from a divisor and detect its image after a
flat cyclotomic base change and inversion of three. The symbol
sketch in \cite[Remark 5.3, footnote 3]{Feld} gives a related use of
symbols. Here we arrange the integral lift on the divisor itself,
so support proves integral generic vanishing before the norm test
is applied. No result from Feld's or Druzhinin's preprint is an
input to this proof.

The companion paper \emph{An integral counterexample to Gersten's
conjecture} \cite[Theorem 1.1]{IntegralDegreeFive} gives an
integral degree-five counterexample at
$A_5=(V_0[x,y]/(5+x^4+y^4))_{(5,x,y)}$, where $V_0$ is the ring of
integers of the unramified degree-eight extension of $\Q_5$.
Its integral cyclotomic lift and Picard-valued specialization test
form a different proof route. Here the quintic model supports a
Steinberg lift and a norm valuation test, proved entirely within
this paper. The two results concern different rings and do not
determine the smallest degree in which integral Gersten injectivity
can fail.

\subsection*{The route through the proof}

Write $C=\X_L$ for the projective generic curve. On the chart $Z\ne0$
let $x=X/Z$ and $y=Y/Z$. The divisor $y=0$ in $\Spec A$ has ring
$S=A/yA$. Section~\ref{sec:model} proves that $S$ is a DVR with
uniformizer $t=x\bmod y$, and that its fraction field $D$ is a totally
ramified degree-five extension of $L$ satisfying
\[
 N_{D/L}(t)=\pi,\qquad (1-t)^5=\zeta.
\]
For a finite extension $E/L$, write $k_E$ for the residue field of
its valuation ring; its residue degree is $[k_E:\F_5]$.
The corresponding closed point $q_0\in C$ has field
$\kappa(q_0)=D$ and specializes to $m$, with
$[D:L]=5$ but $[k_D:\F_5]=1$. Every closed point $q$ specializing
elsewhere satisfies $5\mid[k_{\kappa(q)}:\F_5]$.
This difference in residue degrees supplies the nonvanishing test.

Write $[u]$ for the integral $K_1$ class of a nonzero field element.
Section~\ref{sec:integral} chooses a class
$\beta\in K_2(L;\Z/5)$ whose coefficient boundary in $K_1(L)$ is
$[\zeta]$ and
constructs an integral lift of $a_D=\beta_D\cdot[t]$ in two steps,
where $\beta_D$ is the restriction of $\beta$ to $D$.
The identity $(1-t)^5=\zeta$ turns the coefficient boundary of
$a_D$ into a multiple of the Steinberg symbol
$\{1-t,t\}=0$ in $K_2(D)$ \cite{Matsumoto}. Hence $a_D$ lifts
to an integral class in $K_3(D)$. Quillen's finite-field
calculation gives $K_2(\F_5)=0$ \cite{QuillenFinite}, so localization
then lifts this integral class to $K_3(S)$. Pushforward from $S$
produces $c\in K_3(A)$, and restriction to $\Frac A$ kills it
because its support is the divisor $y=0$.

To detect this candidate, Section~\ref{sec:field} proves that for
every finite extension $E/L$, multiplication by the restriction
$\beta_E$ identifies $E^\times/E^{\times5}$ with
$K_3(E;\Z/5)$. This calculation combines the multiplicative
motivic-to-$K$ spectral sequence of Friedlander and Suslin,
Suslin's product-preserving comparison, and the norm-residue comparison
\cite{FriedlanderSuslin,Suslin,Voevodsky}. The projection formula
identifies field transfer with norm. Normalized valuation modulo
five consequently kills every transferred coefficient class from
a field of residue degree divisible by five, whereas $a_D$ has
transferred valuation one because $N_{D/L}(t)=\pi$.

Section~\ref{sec:detection} proves that this integral class is
nonzero. The scheme $T=\Spec A[1/\pi]$ is an inverse limit of
affine opens of $C$ with flat transition maps. These opens contain
every closed point specializing to $m$, so their omitted points have
residue degree divisible by five. If the coefficient
reduction of $c|_T$ vanished, continuity of coherent $K$-theory
would make it vanish on one such open. Localization
would then express the supported coefficient class
$(q_0)_*(\beta_D\cdot[t])$ as a sum of classes supported at omitted
points. Proper pushforward to $L$ and the norm valuation
give zero for every term on that side and one for $q_0$, a contradiction.
We work in coherent $G$-theory on $C$, which need not be smooth.

\section{The model and its distinguished divisor}\label{sec:model}

We verify the local ring in Theorem~\ref{thm:main} and construct the
divisor on which the integral class will be lifted. The special fiber
forces a factor of five in residue degrees away from $m$, while the
divisor through $m$ has residue degree one and a uniformizer of norm
$\pi$. These properties will distinguish its class from classes
whose supporting points specialize away from $m$. For a finite
extension $E/L$, write $\OO_E$ for its valuation ring and $k_E$
for its residue field.

\begin{lemma}\label{lem:model}
The generic fiber $C=\X_L$ is an integral projective curve. The
ring $A=\OO_{\X,m}$ is a two-dimensional Noetherian regular local
domain of mixed characteristic $(0,5)$ with
$\pi\in\mathfrak m_A^2$. For every closed point $q\in C$
with residue field $E=\kappa(q)$, the $E$-point extends to
$\Spec\OO_E\to\X$. If its closed image is not $m$, then
$5\mid[k_E:\F_5]$.

On the chart $Z\ne0$, set $x=X/Z$, $y=Y/Z$, and let $t$ be the
image of $x$ in $S=A/yA$. Then
\begin{equation}\label{eq:divisor}
 S\cong V[t]/\bigl((t-1)^5+\zeta\bigr)
\end{equation}
is a DVR with uniformizer $t$. Its fraction field $D$ is totally
ramified of degree five over $L$, and
\begin{equation}\label{eq:norm-t}
 N_{D/L}(t)=\pi,\qquad \eta=1-t,\qquad \eta^5=\zeta.
\end{equation}
The equation $Y=0$ cuts out a single closed point $q_0\in C$
with residue field $D$.
\end{lemma}

\begin{proof}
The polynomial
\[
 \Phi_5(1-T)=T^4-5T^3+10T^2-10T+5
\]
is Eisenstein at five. Therefore $\pi=1-\zeta$ is a uniformizer
of $V$, its residue field is $\F_5$, and $5=u\pi^4$ for some
$u\in V^\times$.
The homogeneous equation $P$ defining $\X$ reduces to
\[
 \overline P=X^5-XY^4+Y^5.
\]
If $\alpha$ is a root of $s^5-s+1$, then Frobenius sends
$\alpha$ to $\alpha-1$. Its orbit has length exactly five:
the $r$th iterate is $\alpha-r$, which first returns at $r=5$.
Thus $s^5-s+1$ is irreducible over $\F_5$, and so is its
homogenization $\overline P$. This remains irreducible on adjoining
the independent variable $Z$. A factorization of the primitive
homogeneous polynomial $P$ over $L$ would, by Gauss' lemma, give
primitive positive-degree homogeneous factors over $V$ and hence
a factorization of $\overline P$. Consequently $C$ is integral.
It is a projective plane curve, so it has dimension one.

Put
\[
 F=(x-1)^5-xy^4+y^5+\zeta,\qquad
 R=V[x,y]/(F),\qquad \mathfrak m=(\pi,x,y)\subset R.
\]
In the regular local ring $V[x,y]_{(\pi,x,y)}$ of dimension
three, $F\equiv-\pi$ modulo the square of the maximal ideal.
Thus quotienting by $F$ gives the regular local ring
$A=R_{\mathfrak m}$ of dimension two. A regular local ring is
a domain. Since $F$ is monic in $x$, the ring $R$ is free over
$V[y]$ and hence flat over $V$. Its localization $A$ is also
$V$-flat, so $V$ embeds into $A$. Together with
$A/\mathfrak m_A=\F_5$, this proves mixed characteristic $(0,5)$.
Expanding $F=0$ gives
\[
 \pi=x^5-5x^4+10x^3-10x^2+5x-xy^4+y^5.
\]
In particular, $\pi\in(x,y)A$ and $\mathfrak m_A=(x,y)A$.
Since $5$ is divisible by $\pi$, the term $5x$ is in
$(x,y)^2A$, as are all remaining terms on the right. Hence
$\pi\in\mathfrak m_A^2$. Since $5=u\pi^4$, also
$5\in\mathfrak m_A^2$, so $A$ is ramified in the usual
mixed-characteristic sense. The reduction of the equation has
no linear term at $m$, so the one-dimensional special fiber is
singular there; $A$ is not smooth over $V$.

Modulo $y$, the equation becomes
\[
 f(t)=(t-1)^5+\zeta
 =t^5-5t^4+10t^3-10t^2+5t-\pi.
\]
This is Eisenstein at $\pi$. The finite free $V$-algebra
$B=V[t]/(f)$ is a domain. Every maximal ideal lies over the
maximal ideal of $V$, and
$B/\pi B=\F_5[t]/(t^5)$, so $B$ is local. Since $f(t)=0$
implies $\pi\in(t)B$, its maximal ideal is $(t)$.
Its dimension is one, making it a DVR. The localization defining
$A/yA$ leaves this local ring unchanged, proving
\eqref{eq:divisor}. Its fraction field has degree five over $L$
and residue degree one, hence total ramification. The constant
term gives
\[
 N_{D/L}(t)=(-1)^5f(0)=\pi.
\]
The relation $(1-t)^5=\zeta$ also follows from $f(t)=0$.
As $S[1/\pi]=D$, the equation $y=0$ on the generic affine
chart gives a single closed point. There is no additional point
with $Y=Z=0$ on $C$, since substitution in $P$ would force
$X=0$. This proves the assertion about $q_0$ on all of $C$.

We have constructed the distinguished point with residue degree one.
It remains to establish the residue-degree constraint on points that
specialize away from $m$.
A closed point of $C$ has residue field finite over $L$.
Properness of $\X/V$ extends its $E$-point to $\Spec\OO_E$.
On the special fiber $Y=0$ forces $X=0$, so $m$ is the only
point with $Y=0$. If the closed image is elsewhere, $X/Y$ in
$k_E$ is a root of $s^5-s+1$. Thus $k_E$ contains $\F_{5^5}$,
and $5\mid[k_E:\F_5]$.
\end{proof}

\section{Lifting on the divisor}\label{sec:integral}

The divisor $\Spec S\subset\Spec A$ now supplies an integral class
with zero generic restriction. We first lift a coefficient class
to $K_3(D)$ and then extend that lift across the closed point of $S$.
Its coefficient image will be tested for nonvanishing only after
this construction.

Let $M_5$ be the mod-five Moore spectrum. We use
$K_i(B;\Z/5)=\pi_i(K(B)\wedge M_5)$, and the same convention for
$G$-theory. The coefficient sequence is
\begin{equation}\label{eq:coeff}
 K_i(B)\xrightarrow{5}K_i(B)\xrightarrow{\rho_B}K_i(B;\Z/5)
 \xrightarrow{\partial_{\mathrm{coef}}}K_{i-1}(B)
 \xrightarrow{5}K_{i-1}(B).
\end{equation}
Here $\rho_B$ denotes coefficient reduction, and
$\partial_{\mathrm{coef}}$ denotes the connecting map.
Multiplication of a coefficient class by an integral class uses
the natural $K(B)$-module structure on $K(B)\wedge M_5$.

For an element $u$ of a field's multiplicative group, write $[u]$
for its integral $K_1$ class. Since
$[\zeta]\in K_1(L)=L^\times$ has order five, exactness permits
us to fix
\[
 \beta\in K_2(L;\Z/5),\qquad
 \partial_{\mathrm{coef}}\beta=[\zeta].
\]
For every finite extension $E/L$, write $\beta_E$ for its restriction.
The class to be lifted is
\[
 a_D=\beta_D\cdot[t]\in K_3(D;\Z/5).
\]

Resolution identifies $K$-theory with coherent
$G$-theory for the regular rings $S$ and $A$; closed-immersion
pushforward is taken in $G$-theory
\cite[\S7.1 and \S7, paragraph 2.7 and the discussion following (2.8)]{QuillenHigher}.

\begin{proposition}\label{prop:integral}
There exists $b\in K_3(S)$ whose restriction to $D$, followed
by coefficient reduction, is $a_D$. If
$j:\Spec S\hookrightarrow\Spec A$ and $c=j_*b\in G_3(A)=K_3(A)$,
then $c|_{\Frac A}=0$. Moreover, for
$T=\Spec A[1/\pi]$ and $i:\Spec D\hookrightarrow T$,
\[
 c|_T\bmod5=i_*a_D\quad\text{in }G_3(T;\Z/5).
\]
\end{proposition}

\begin{proof}
The coefficient cofiber sequence is a sequence of $K(D)$-modules.
Its connecting map is therefore compatible with multiplication by
an integral $K(D)$ class, up to the suspension sign. By \eqref{eq:norm-t},
$\zeta=\eta^5$ and $t=1-\eta$. Both $\eta$ and $1-\eta$
are nonzero: the first has fifth power $\zeta\ne0$, and
the second is the uniformizer $t$ of $S$. Matsumoto's
Steinberg relation \cite[\S5, Lemma 5.6, Theorem 5.10, and
Corollary 5.11]{Matsumoto}, with the identification of symbols and
products in Quillen $K_2(D)$
\cite[III.6.1, IV.1.7.1, IV.1.10.1, and IV.7.2]{Weibel}, gives
\begin{equation}\label{eq:steinberg}
 \partial_{\mathrm{coef}}a_D
 =\pm[\zeta]\cdot[t]
 =\pm5\{\eta,1-\eta\}=0.
\end{equation}
The sign does not affect the vanishing. Exactness of
\eqref{eq:coeff} yields an integral lift
$\widetilde a_D\in K_3(D)$ of $a_D$.
The boundary $[\zeta]\in K_1(D)[5]$ of $\beta_D$ itself is
still nonzero, even though $\zeta=\eta^5$. Being a fifth power
makes the class of $\zeta$ vanish in $D^\times/D^{\times5}$,
not in the integral group $K_1(D)=D^\times$.
The Steinberg relation is what kills the boundary of the product
$a_D$.

Localization and d\'evissage for the DVR $S$ give
\cite[\S7, Proposition 3.2]{QuillenHigher}
\[
 K_3(S)\longrightarrow K_3(D)\longrightarrow K_2(\F_5).
\]
The last group is zero by Quillen's finite-field computation
\cite{QuillenFinite}. We may therefore choose
$b\in K_3(S)$ restricting to $\widetilde a_D$. This is a
second lift, distinct from the coefficient lift to $K_3(D)$.

Let $c=j_*b\in G_3(A)=K_3(A)$. Its restriction to
$\Spec A[1/y]$ is zero because the closed immersion $j$
has empty pullback there. The element $y$ is nonzero in the
domain $A$: its quotient $S$ has dimension one, whereas $A$
has dimension two. Hence $y$ becomes invertible in $\Frac A$,
and $c|_{\Frac A}=0$ integrally.

Base change of $j$ along $T\to\Spec A$ is
$i:\Spec D\hookrightarrow T$, since $S[1/\pi]=D$.
Flat base change for coherent pushforward
\cite[V.3.7.2]{Weibel} here is the natural isomorphism
\[
 A[1/\pi]\otimes_A M\cong D\otimes_S M
\]
for a finite $S$-module $M$, with both sides viewed as
$A[1/\pi]$-modules. Thus restriction of scalars commutes with
localization. Naturality of coefficient reduction gives the
commutative diagram
\[
\begin{array}{ccc}
 K_3(S)&\xrightarrow{\ j_*\ }&G_3(A)\\
 \big\downarrow&&\big\downarrow\\
 K_3(D;\Z/5)&\xrightarrow{\ i_*\ }&G_3(T;\Z/5).
\end{array}
\]
Both vertical arrows mean restriction followed by coefficient
reduction. The left arrow sends $b$ to $a_D$,
which proves the stated identity.
\end{proof}

\section{The field comparison and norm detector}\label{sec:field}

Proposition~\ref{prop:integral} constructs an integral class $c$ with
zero generic restriction. To prove it nonzero, we will test its
coefficient image by proper pushforward to $L$ and a valuation.
Localization on the projective curve will introduce arbitrary classes
in $K_3(E;\Z/5)$, for finite extensions $E/L$ arising as
closed-point fields.
We therefore need a description of all these groups that identifies
field transfer with norm.

Keep the class $\beta$ and its restrictions $\beta_E$ fixed as in
Section~\ref{sec:integral}. For every finite extension $E/L$, let
$k_E$ be its residue field, and normalize $v_E$ so that a uniformizer
has valuation one.

\begin{lemma}[Multiplicative field comparison]\label{lem:field}
For $L=\Q_5(\zeta)$ and the fixed $\beta$ above, every finite
extension $E/L$ has an isomorphism
\begin{equation}\label{eq:theta}
 \theta_E:E^\times/E^{\times5}\xrightarrow{\sim}K_3(E;\Z/5),
 \qquad u\bmod E^{\times5}\longmapsto\beta_E\cdot[u].
\end{equation}
The degree-three motivic edge identifies the target with
$H^1_{\mathrm{et}}(E,\mu_5^{\otimes2})$. Under this edge,
the displayed map is Kummer reduction followed by cup product
with the nonzero weight-one edge image
$\alpha_E\in H^0(E,\mu_5)$ of $\beta_E$.
\end{lemma}

\begin{proof}
We first identify the degree-three group, then locate the image of
$\beta_E$ in degree two, and finally compute multiplication by
$\beta_E$ on integral $K_1$. The strongly convergent motivic-to-$K$
spectral sequence for a characteristic-zero field is
\begin{equation}\label{eq:motivic}
 E_2^{a,b}=H^{a-b}(E,\Z/5(-b))
 \quad\Longrightarrow\quad K_{-a-b}(E;\Z/5).
\end{equation}
It has a pairing with the integral spectral sequence that induces
the integral/coefficient $K$-module product on the abutment
\cite[Theorem 16.2 and Appendix C.2--C.3]{FriedlanderSuslin}.
The comparison with motivic complexes preserves products
\cite[Lemma 3.4 and Theorem 6.1]{Suslin}.
Write $s=a-b$ and $j=-b$, so a term $H^s(E,\Z/5(j))$ has total
$K$-degree $2j-s$. Negative weights vanish. For $j\geq0$, the
field motivic groups vanish for $s>j$, and the norm-residue and
motivic--\'etale comparison theorems identify them, for every
$s\leq j$, with
\[
 H^s_{\mathrm{et}}(E,\mu_5^{\otimes j});
\]
see \cite[Theorems 6.16--6.17]{Voevodsky} and
\cite[VI.4.1--4.2]{Weibel}. In particular, the groups vanish for
$s<0$. These comparisons apply because $5$ is invertible in $E$.

The $5$-cohomological dimension of a finite extension of $\Q_5$
is two \cite[VI.\S7, before Proposition 7.3]{Weibel}. Thus the
only potentially nonzero terms satisfy $0\leq s\leq\min(j,2)$.
For $r\geq2$ a differential changes $(s,j)$ to
$(s+2r-1,j+r-1)$, so none can join two nonzero terms. Total
degree three has the single term $(s,j)=(1,2)$, yielding the
natural identification
\begin{equation}\label{eq:degree-three-edge}
 K_3(E;\Z/5)\simeq H^1_{\mathrm{et}}(E,\mu_5^{\otimes2}).
\end{equation}
In particular, this $K$-group is killed by five. In total degree
two the weight-one quotient of the filtration is $H^0(E,\mu_5)$;
the only other possible piece is $H^2(E,\mu_5^{\otimes2})$ in
weight two.

Let $\alpha_E$ be the weight-one image of $\beta_E$.
To see that it is nonzero, restrict to an algebraic closure
$\overline E$. The same comparison there has only $s=0$ terms,
so its total degree-two group has just the weight-one piece
$H^0(\overline E,\mu_5)$. Matsumoto's presentation of field
$K_2$ (\cite[\S5]{Matsumoto}; see also \cite[III.6.1]{Weibel}) gives
$K_2(\overline E)/5=0$: every symbol
$\{u,v\}$ equals $5\{r,v\}$ after choosing $r^5=u$.
The coefficient sequence therefore identifies
$K_2(\overline E;\Z/5)$ with $K_1(\overline E)[5]$ through
$\partial_{\mathrm{coef}}$. The restriction of $\beta_E$ is
nonzero, since its boundary is the primitive root $[\zeta]$.
Naturality of the filtration shows that $\alpha_E$ has nonzero
image in $H^0(\overline E,\mu_5)$. It is consequently a generator
of $H^0(E,\mu_5)$; only this nonvanishing is needed.

It remains to determine the actual product from these filtration
pieces. Integral $K_1(E)=E^\times$ has only the weight-one piece
$H^1(E,\Z(1))$: on the diagonal $2j-s=1$, weights $j\geq2$
give $s>j$ and vanish, weight zero gives
$H^{-1}(E,\Z(0))=0$, and negative weights vanish. The
identification $\Z(1)\simeq\mathbb G_m[-1]$ identifies reduction
of $[u]$ with its Kummer class in
$H^1(E,\mu_5)=E^\times/E^{\times5}$.

Write $F^j$ for the decreasing weight filtration. The preceding
calculations give
\[
 F^1K_1(E)=K_1(E),\quad F^2K_1(E)=0,\qquad
 F^2K_3(E;\Z/5)=K_3(E;\Z/5),\quad F^3K_3(E;\Z/5)=0.
\]
The class $\beta_E$ lies in $F^1K_2(E;\Z/5)$ and has image
$\alpha_E$ in its quotient by $F^2$. The mixed filtered pairing
satisfies
\[
 F^2K_2(E;\Z/5)\cdot F^1K_1(E)
 \subseteq F^3K_3(E;\Z/5)=0.
\]
Thus the possible weight-two part of $\beta_E$ contributes nothing:
the product
$\beta_E\cdot[u]$ is determined by $\alpha_E$ and the
weight-one class of $[u]$. Product compatibility identifies it,
under \eqref{eq:degree-three-edge}, with
$\alpha_E\cup[u]$.
Since $\alpha_E$ generates the constant one-dimensional
$\F_5$-module $\mu_5$, this cup product is an isomorphism.
The product factors through fifth powers because its target is
killed by five and $[u^5]=5[u]$. This proves
\eqref{eq:theta} with its asserted multiplicative identification.
\end{proof}

In particular, the coefficient class already constructed on the
divisor is $a_D=\theta_D(t)$. We now relate the comparison maps for
different extensions of $L$.

\begin{lemma}[Norm detector]\label{lem:norm-detector}
For every finite extension $E/L$
and every $u\in E^\times$,
\begin{equation}\label{eq:transfer}
 \Tr_{E/L}\theta_E(u)=\theta_L(N_{E/L}u).
\end{equation}
Consequently
\[
 \lambda_L=(v_L\bmod5)\circ\theta_L^{-1}:
 K_3(L;\Z/5)\longrightarrow\Z/5
\]
satisfies
\[
 \lambda_L\bigl(\Tr_{E/L}(a)\bigr)=0
 \quad\text{for every }a\in K_3(E;\Z/5)
 \quad\text{if }5\mid[k_E:k_L].
\]
\end{lemma}

\begin{proof}
Restriction of scalars obeys the tensor projection formula at the
level of $K$-theory spectra \cite[V.3.3.2]{Weibel}. Smashing its
base-field factor with $M_5$ gives the mixed pairing identity
\[
 \Tr_{E/L}(\beta_E\cdot[u])
 =\beta\cdot\Tr_{E/L}[u]
 =\beta\cdot[N_{E/L}u].
\]
The second equality uses the norm description of field $K_1$
transfer \cite[III.1.7.1 and the following paragraph]{Weibel}. This proves
\eqref{eq:transfer}.

For normalized valuations, the norm of the prime ideal of $E$
is the $[k_E:k_L]$th power of the prime ideal of $L$. Hence
\[
 v_L(N_{E/L}u)=[k_E:k_L]v_E(u).
\]
Every $a\in K_3(E;\Z/5)$ has the form $\theta_E(u)$ by
Lemma~\ref{lem:field}. Applying the formula to that $u$ proves
the assertion for all $a$. The factor is the residue degree,
not the total degree $[E:L]$.
\end{proof}

\section{Detection on the projective curve}\label{sec:detection}

Proposition~\ref{prop:integral} has constructed an integral class $c$
with $c|_{\Frac A}=0$ and coefficient image $i_*a_D$ on
\[
 T=\Spec A[1/\pi].
\]
By Lemma~\ref{lem:field}, $a_D=\beta_D\cdot[t]=\theta_D(t)$.
To prove its pushforward $i_*a_D$ nonzero, we extend the supported
class to the projective curve $C$, where proper pushforward permits
the norm test. The scheme $T$ is an inverse limit of
affine opens of $C$ with flat transition maps, and need not itself
be an open subset. The following lemma makes that passage explicit.

\begin{lemma}\label{lem:opens}
Let $\Sigma=R\setminus\mathfrak m$, where
$R=V[x,y]/(F)$ and $\mathfrak m=(\pi,x,y)$.  For $h\in\Sigma$ put
\[
 U_h=\Spec R[1/\pi,1/h]\subset C.
\]
Then the refinement maps $U_{hk}\to U_h$ are affine and flat, and
\begin{equation}\label{eq:continuity}
 T=\varprojlim_{h\in\Sigma}U_h,
 \qquad
 \varinjlim_{h\in\Sigma}G_3(U_h;\Z/5)
   \xrightarrow{\ \sim\ }G_3(T;\Z/5).
\end{equation}
Every $U_h$ contains $q_0$.  Its complement in $C$ is a finite set of
closed points $q$, and their residue fields $E_q=\kappa(q)$ satisfy
$5\mid[k_{E_q}:\F_5]$.
\end{lemma}

\begin{proof}
The multiplicative set $\Sigma$ is directed by taking products:
$hk$ gives a common refinement of $h$ and $k$.  Transitivity of
localization gives
\[
 \varinjlim_{h\in\Sigma}R[1/\pi,1/h]
 =R_{\mathfrak m}[1/\pi]=A[1/\pi].
\]
Each refinement is a principal localization, so the corresponding
scheme map is affine and flat.  The stages and their limit are affine
Noetherian schemes, hence separated.  Quillen's continuity theorem,
in its coherent $K'$ conclusion
\cite[\S7, Proposition~2.2]{QuillenHigher}, therefore applies to
their integral $G$-groups.

To pass to coefficients, use the natural short exact sequence
\[
 0\longrightarrow G_3(Y)/5G_3(Y)
 \longrightarrow G_3(Y;\Z/5)
 \longrightarrow G_2(Y)[5]\longrightarrow0
\]
for each stage $Y=U_h$ and for $Y=T$.  Filtered colimits of abelian
groups are exact, and thus commute with the quotient by five and
the kernel of multiplication by five.  Integral continuity in
degrees three and two identifies the outer terms in the resulting
comparison of short exact sequences.  It consequently identifies
the middle terms, proving the coefficient assertion
in~\eqref{eq:continuity}.

Since $h\notin\mathfrak m$, its image in the local divisor ring $S$
is a unit.  Thus $q_0=\Spec D$ belongs to $U_h$.  More generally,
consider a closed point $q\in C$ with residue field $E$ whose
extension $\Spec\OO_E\to\X$ specializes to $m$.
The inverse image of the chart $Z\ne0$ contains the closed point of
$\Spec\OO_E$, and hence is all of $\Spec\OO_E$.
On that chart the reduction of $h$ at $m$ is nonzero.  Its pullback
to $\OO_E$ is therefore a unit, so $q$ lies in $U_h$.
It follows that each point omitted from $U_h$ specializes away
from $m$.  Lemma~\ref{lem:model} then gives
$5\mid[k_E:\F_5]$.  Finally, $C$ is an integral projective curve
and $U_h$ is nonempty, so its closed complement consists of
finitely many closed points.
\end{proof}

\begin{proof}[Proof of Theorem~\ref{thm:main}]
Take the integral class $c\in K_3(A)$ constructed in
Proposition~\ref{prop:integral}, and put
\[
 w=(q_0)_*\theta_D(t)\in G_3(C;\Z/5).
\]
Coherent $G$-theory provides proper pushforward even when $C$ is singular.
The point immersion is finite, and the structural map
$p:C\to\Spec L$ is proper with $C$ carrying the ample line bundle
$\mathcal O_C(1)$.  Coherent pushforward and its composition law
therefore apply to both maps
\cite[\S7, paragraph 2.7 and the discussion following (2.8)]{QuillenHigher}.

Flat base change for the point immersion on each $U_h$
\cite[V.3.7.2]{Weibel}, followed by~\eqref{eq:continuity}, identifies
$w|_T$ with the pushforward of $\theta_D(t)$ along
$\Spec D\hookrightarrow T$.
By Proposition~\ref{prop:integral}, this is precisely the coefficient
reduction of $c|_T$.

Suppose that $w|_T=0$.  The element $w|_{U_1}$ represents this class
in the filtered colimit~\eqref{eq:continuity}.  A zero element in
that colimit becomes zero at a finite stage, so
$w|_{U_h}=0$ for some $h\in\Sigma$.
Give $Z_h=C\setminus U_h$ its reduced closed structure.
Apply Moore coefficients to coherent localization and d\'evissage
at these closed points \cite[\S7, Proposition~3.2 and
paragraph~3.4]{QuillenHigher}. The resulting exact segment is
\[
 \bigoplus_{q\in Z_h}K_3(E_q;\Z/5)
 \longrightarrow G_3(C;\Z/5)
 \longrightarrow G_3(U_h;\Z/5).
\]
Consequently there are classes $a_q\in K_3(E_q;\Z/5)$ such that
\begin{equation}\label{eq:omitted-relation}
 w=\sum_{q\in Z_h}q_*a_q.
\end{equation}
These are arbitrary coefficient classes supplied by localization.
Lemma~\ref{lem:field} applies to every finite extension $E_q/L$,
so its surjectivity writes each of them as
$a_q=\theta_{E_q}(u_q)$ for some $u_q\in E_q^\times$.

Apply the proper pushforward
$p_*:G_3(C;\Z/5)\to K_3(L;\Z/5)$ to
\eqref{eq:omitted-relation}.
For every closed point, the composite $p_*q_*$ is the field transfer.
Lemma~\ref{lem:opens} gives $5\mid[k_{E_q}:\F_5]$, and
Lemma~\ref{lem:norm-detector} therefore gives
\[
 \lambda_L\bigl(p_*q_*a_q\bigr)
 =v_L\bigl(N_{E_q/L}(u_q)\bigr)\bmod5
 =[k_{E_q}:\F_5]v_{E_q}(u_q)\bmod5=0.
\]
Thus \eqref{eq:omitted-relation} forces $\lambda_L(p_*w)=0$.
On the other hand, the distinguished point has residue field $D$,
and its parameter satisfies $N_{D/L}(t)=\pi$.  The same transfer
formula yields
\[
 \lambda_L(p_*w)
 =\lambda_L\bigl(\operatorname{Tr}_{D/L}\theta_D(t)\bigr)
 =v_L\bigl(N_{D/L}(t)\bigr)\bmod5
 =v_L(\pi)\bmod5=1.
\]
This contradiction proves $w|_T\ne0$.
Hence the coefficient reduction of $c|_T$ is nonzero, so the
integral class $c\in K_3(A)$ is itself nonzero.
Its image at the fraction field is zero by
Proposition~\ref{prop:integral}; together with the properties of $A$
in Lemma~\ref{lem:model}, this proves the theorem.
\end{proof}

\bibliographystyle{plain}
\bibliography{refs/sources}
\end{document}
